% RESEARCH_DRAFT_SHA256: 35e078bb95bfd23f17374c71284ab6f26029f757295c0786ea51a0f6a5b74625
% Layout and source-reference wrappers expanded for the web renderer.
\documentclass[11pt,a4paper]{article}
\usepackage[margin=25mm]{geometry}
\usepackage{fontspec}
\setmainfont{lmroman10-regular.otf}[BoldFont=lmroman10-bold.otf,
 ItalicFont=lmroman10-italic.otf,BoldItalicFont=lmroman10-bolditalic.otf]
\usepackage{amsmath,amssymb,mathtools}
\usepackage{unicode-math}
\setmathfont{latinmodern-math.otf}
\AtBeginDocument{\let\setminus\smallsetminus}
\usepackage{xurl,hyperref}
\hypersetup{colorlinks=true,urlcolor=blue}
\setcounter{secnumdepth}{0}
\setlength{\parindent}{0pt}
\setlength{\parskip}{5pt}
\setlength{\emergencystretch}{3em}
\begin{document}
\section*{Research-attempt provenance}
This standalone research attempt retains the complete original definition body below,
including its historical source and preparation notes. Those notes describe the
earlier catalogue preparation; the new attempt and its verification are identified
separately after that body. No original ID, catalogue status, or source statement is changed.

\section*{Dynamic Planar Nearest Neighbors}\label{5500063}
{\small UnsolvedMath ID 5500063 · AMR-054-0063 · Computer Science · AMR Open Problem Lists}
\subsection{Definitions and mathematical statement}
Maintain a finite set $P\subset\mathbb R^2$ under insertions of
new sites and deletions of existing sites. A query supplies
$q\in\mathbb R^2$ and must return a site $p\in P$ minimizing
$\|q-p\|_2$. Any minimizer is acceptable in a tie; when $P$ is
empty, the structure reports this fact. Distances are exact,
not approximate.

Is there a data structure supporting insertions, deletions,
and nearest-neighbor queries all in $O(\log(n+2))$ time, where
$n$ is the current number of sites? Use the usual real-arithmetic
model for planar geometric data structures, with polynomial
space and preprocessing. The strongest reading requests a
worst-case bound for every operation. An amortized or randomized
solution must state its guarantee explicitly; these are distinct
performance formulations of the source's general target.

Updates may occur in arbitrary order, and their locations are
not assumed to follow a random distribution. A structure allowing
only insertions does not meet the fully dynamic requirement.
Nor does logarithmic query time alone suffice when one update
operation is slower. For an amortized formulation, the running
time is assessed over whole update sequences, rather than
charging every deletion its own worst-case cost. The source's
related convex-hull reduction can guide constructions but does
not itself supply the requested simultaneous bounds.
\subsection{Short English statement}
Can a fully dynamic planar point set support exact nearest-neighbor queries, insertions, and deletions in logarithmic time per operation?
\subsection{Sources}
\begin{itemize}
\item[S1] ULAM.AI and the UnsolvedMath Contributors, supplied problems.json, record 5500063 (AMR-054-0063), AMR Open Problem Lists. Catalogue identity and source transcription; CC BY 4.0.
\url{https://huggingface.co/datasets/ulamai/UnsolvedMath}.
\item[S2] Open Problems Project - Computational Geometry, Problem 63. Original source indexed by AMR-054-0063.
\url{https://topp.openproblem.net/p63}.
\item[S3] The Open Problems Project, original numbered problem and status notes.
\url{https://topp.openproblem.net/p63}.
\item[S4] Incremental Planar Nearest Neighbor Queries with Optimal Query Time, SoCG2025; insertion-only progress and comparison with fully dynamic bounds.
\url{https://drops.dagstuhl.de/storage/00lipics/lipics-vol332-socg2025/html/LIPIcs.SoCG.2025.59/LIPIcs.SoCG.2025.59.html}.
\end{itemize}

\section{Literature and context check}
TOPP Problem 63 was consulted for the fully dynamic target and its
three-dimensional convex-hull reduction. The argument below tests explicit
Voronoi maintenance, not every possible implicit data structure. No insertion-only
result is promoted to the fully dynamic worst-case guarantee.
\section{Research attempt: Why explicit Voronoi maintenance cannot charge every update logarithmically}
The usual lifting identity is immediate: for $p=(a,b)$ and $q=(u,v)$,
\[
 \|q-p\|^2=\|q\|^2+(-2u,-2v,1)\cdot(a,b,a^2+b^2).
\]
Thus a minimum linear-functional query on the lifted points returns a nearest
site. Insertions and deletions transfer one-for-one. This is an exact reduction,
but supplies no data structure by itself.

The attempted planar implementation stored every Voronoi adjacency explicitly
and updated only a supposedly bounded neighborhood after an insertion. Let
$p_1,\ldots,p_m$ be distinct sites on the unit circle, and insert the origin.
For each $i$, at $q_i=p_i/2$ the origin and $p_i$ are both at distance $1/2$.
For $j\ne i$,
\[
 \|q_i-p_j\|^2=\tfrac54-p_i\cdot p_j>\tfrac14.
\]
There is a strict positive gap to every other site. By continuity, a small
segment of the origin--$p_i$ perpendicular bisector about $q_i$ belongs to
their shared Voronoi boundary. Hence the inserted site has $m$ distinct new
Voronoi adjacencies. A representation allocating a separate record for every
new adjacency must write at least $m$ records on this one insertion. Deleting
the origin reverses the example.

The initial circle creates a high-order Voronoi vertex, but the adjacency
conclusion does not rely on retaining that degeneracy. For fixed finite $m$,
the finitely many displayed strict gaps have a positive minimum. A sufficiently
small perturbation of the sites preserves positivity at suitably adjusted
bisector points. In particular the effect cannot be dismissed as a tie-handling
artifact. The original problem in any event permits the circular input.

This defeats the proposed local-change bound for an explicit diagram, not the
existence of a more economical representation. A lazy representation might
describe the new cell without eagerly materializing all its adjacencies; then
one must analyze exact queries and future deletions using that representation.
Similarly dividing sites into binary-sized static blocks gives at most
$\sum_{k\le\log_2n}O(k+1)=O(\log^2(n+2))$ query work if each nonempty block
is queried separately in logarithmic time. This upper-bound calculation does
not become $O(\log n)$ merely because there are logarithmically many blocks.
\section{Verification and adversarial audit}
The circle inequality was checked for every pair $i\ne j$, including
antipodal points. The $m$ new adjacencies are certified by separate strict local
distance comparisons. Empty sets and ties remain governed by the source
convention. The write-count argument is expressly representation-specific.
\paragraph{Audit corrections.} The linear explicit-diagram update was not stated as a lower bound for implicit structures, and the lifting reduction was not mistaken for an algorithm.
\paragraph{External computation.} None; symbolic arguments and boundary checks only.
\section{Scope and remaining gaps}
An implicit fully dynamic structure with all three logarithmic guarantees
has not been constructed. The first gap is handling updates with many geometric
changes without paying for all changes, while keeping exact worst-case queries
and arbitrary later deletions efficient.
\paragraph{Integration review.} Codex (GPT-6; exact serving revision not exposed) reviewed the written arguments and expanded layout and source-reference wrappers. Repository metadata and the production web renderer were checked using Python 3.13 and JavaScript with MathJax 3.2.2. This model-assisted review is not independent human proof verification; no PDF compilation or formal proof checking is asserted.

\section{Completion Estimate}
\noindent\textbf{COMPLETION: 0\%}
The calculations below the preserved source statement investigate a concrete route,
but establish no substantiated new progress toward the entire original target.
Known special cases and failed sufficient conditions are not credited as a solution.
\section{References}
The original source references and catalogue attribution are preserved in the
Sources section above. The following URLs were actually consulted for this attempt;
retention of an additional original reference does not claim fresh consultation of it.
\begin{itemize}
\item \url{https://topp.openproblem.net/p63} (consulted 2026-10-11). 
\end{itemize}
\section{Model and contributor disclosure}
OpenAI gpt-6-astra, requested reasoning configuration ultra, generated this research
attempt on 2026-10-11. The exact serving revision is not exposed. Codex GPT-6
performed orchestration and repository-format review; its exact serving revision
is also not exposed. The mathematical route, boundary checks and final adversarial
self-audit followed contribution prompts A then B in
\texttt{docs/contribution\_prompts.md}, after reading \texttt{CONTRIBUTING.md}.
Python 3.14.2 assembled the standalone documents; any mathematical computation is
specified in Verification. No theorem prover was used. No human mathematical
authorship, substantial human proof correction, or independent human proof checking
is asserted. Model-assisted review, source retrieval, and structural checks do not
establish correctness or novelty. Full private reasoning transcripts are unavailable
as reproducibility artifacts; the complete arguments, inputs and outcomes actually
used are stated here or in the companion batch audit.
\section{Final status}
Unresolved. This is a research attempt with an explicit failed route or restricted
calculation, not a full solution, not a current-status certification, and not a claim
that the elementary material is new. The catalogue's original source-status label
is preserved independently of this attempt status.
\end{document}
